
\medskip
\phantomsection\label{lem:sga3-VIB-11.17.0}%
\noindent\textbf{Lemma 11.17.0.---}%
\footnote{\label{note:sga3-VIB-130}Editor's note: we have
added this lemma, drawn from the proof of \SGARefLink{par:sga3-VIB-5.6}{Theorem 5.6} of
\textup{[DG70]}, Section III.3. By an abuse of notation, the
same letter \(E\) denotes a \(k\)-vector space and the
\(k\)-module scheme
\(W(E)=\operatorname{Spec}S(E^*)\).}
\emph{Let \(k\) be an algebraically closed field, let \(G\)
be a reduced affine algebraic \(k\)-group, let \(V\) be a
finite-dimensional \(G\)-module over \(k\), and let
\(v\in V\). Let \(E\) be the vector subspace of \(V\)
spanned by the vectors \(gv\), for \(g\in G(k)\). Then
\(E\) is the smallest \(G\)-submodule of \(V\) containing
\(v\). Consequently, the morphism}
\[
  G\times E\longrightarrow V,\qquad (g,x)\longmapsto gx,
\]
\emph{factors through \(E\).}

\emph{Proof.}
By \hyperref[lem:sga3-VIB-11.8]{11.8}, there is a smallest
\(G\)-submodule \(U\) of \(V\) containing \(v\). Indeed,
if
\[
  \mu:V\longrightarrow V\otimes\mathcal A(G)
\]
denotes the comodule structure and
\[
  \mu(v)=\sum_{i=1}^{n}v_i\otimes a_i
\]
with the \(a_i\) linearly independent, then
\[
  U=\operatorname{Span}_k(v_1,\ldots,v_n).
\]
It is clear that \(U\) contains \(E\), and that the
morphism \(G\times U\to V\) factors through \(U\).

Conversely, the inverse image of \(E\) under the morphism
\[
  \mu_v:G\longrightarrow V,\qquad g\longmapsto gv,
\]
is a closed subset of \(G\) containing all the rational
points. Those points are dense in \(G\), since \(G\) is of
finite type over \(k\)
\textup{(cf.\ EGA IV\(_3\), 10.4.8)}. Hence
\[
  \mu_v^{-1}(E)=G.
\]
Since \(G\) is reduced, \(\mu_v\) therefore factors through
\(E\), whence \(E=U\).

\medskip
\phantomsection\label{thm:sga3-VIB-11.17}%
\noindent\textbf{Theorem 11.17 \textup{(Chevalley)}.---}
\emph{Let \(k\) be a field, let \(G\) be an affine
\(k\)-group \textup{(not necessarily of finite type)}, and
let \(N\) be a closed group subscheme of \(G\), invariant
in \(G\). Then the \textup{(fpqc)} sheaf quotient \(G/N\)
is representable by an affine \(k\)-group.}%
\footnote{\label{note:sga3-VIB-131}Editor's note: for
another proof of this theorem that does not use the results
of Expos\'e VI\(_A\), see \textup{[Ta72]}, \SGARefLink{prop:sga3-VIB-5.2}{Theorem 5.2}
\textup{(see also \hyperref[rem:sga3-VIB-11.18.5]{Remark 11.18.5})}.}

\emph{Proof.}
Suppose first that \(G\) is of finite type. By
\hyperref[thm:sga3-VIA-3.2]{VI\(_A\), 3.2} and
\hyperref[par:sga3-VIA-5.2]{5.2}, the \textup{(fpqc)}
sheaf quotient \(G/N\) is representable by a \(k\)-group
\(Q\); it remains to prove that \(Q\) is affine. We proceed
in several steps, first supposing that \(k\) is
algebraically closed.%
\footnote{\label{note:sga3-VIB-132}Editor's note: in what
follows, we have expanded the original and corrected the
erroneous assertion
\((G/N)_{\mathrm{red}}=G_{\mathrm{red}}/N_{\mathrm{red}}\),
drawing on \textup{[DG70]}, Section III.3, \SGARefLink{par:sga3-VIB-5.6}{5.6}.}

\smallskip
\phantomsection\label{case:sga3-VIB-11.17-a}%
\noindent\textup{(a)}
Suppose, in addition, that \(G\) is reduced and connected,
and that \(N\) is reduced. By
\hyperref[thm:sga3-VIB-11.16]{11.16}, there are a
finite-dimensional \(G\)-module \(V\) over \(k\) and a
line \(D=ke_0\) such that
\[
  \operatorname{Norm}_G(D)=N.
\]
In particular, \(N\) acts on \(D\) through a character
\[
  \chi:N\longrightarrow\mathbb G_{m,k}.
\]

Fix \(h\in N(k)\). For every \(g\in G(k)\), one has
\[
  hge_0=g(g^{-1}hg)e_0
       =\chi(g^{-1}hg)\,ge_0.
\]
Thus \(\chi(g^{-1}hg)\) is an eigenvalue of \(h\). The
continuous map
\[
  \varphi:G(k)\longrightarrow k,\qquad
  g\longmapsto\chi(g^{-1}hg),
\]
therefore takes only finitely many values. Since \(G(k)\)
is irreducible \textup{(it is dense in \(G\))}, it follows
that
\[
  \varphi(g)=\varphi(e)=\chi(h)
  \qquad\text{for every }g\in G(k).
\]
Consequently,
\phantomsection\label{form:sga3-VIB-11.17-character-invariance}%
\[
  \chi(g^{-1}hg)=\chi(h)
  \qquad
  \text{for all }g\in G(k),\ h\in N(k).
\]

Let \(E\) be the vector subspace of \(V\) spanned by the
vectors \(ge_0\), for \(g\in G(k)\). By
\hyperref[lem:sga3-VIB-11.17.0]{Lemma 11.17.0}, this is the
\(G\)-submodule of \(V\) generated by \(e_0\).

By the preceding identity, the two morphisms
\[
  N\times E\longrightarrow E,\qquad
  (h,x)\longmapsto hx
  \quad\text{and}\quad
  (h,x)\longmapsto\chi(h)x,
\]
agree on
\[
  (N\times E)(k)=N(k)\times E(k),
\]
which is dense in \(N\times E\). Since \(N\times E\) is
reduced \textup{(and \(E\) is separated)}, the two
morphisms are equal. Thus \(N\) acts on \(E\) by scalar
multiplications.

It follows that \(N\) is contained in the kernel \(K\) of
the morphism
\[
  \rho:G\longrightarrow
    \operatorname{GL}\bigl(\operatorname{End}_k(E)\bigr)
\]
defined, for every \(k\)-algebra \(R\), by
\[
  \rho(g)(u)=gug^{-1}
  \qquad
  \bigl(g\in G(R),\
  u\in\operatorname{End}_R(E\otimes_kR)\bigr).
\]
Conversely, if \(g\in K(R)\), then
\[
  g(Re_0)=Re_0,
\]
and hence \(g\in N(R)\). Thus \(N=K\). By
\hyperref[prop:sga3-VIA-5.4.1]{VI\(_A\), 5.4.1}, the
morphism
\[
  G/N\longrightarrow
  \operatorname{GL}\bigl(\operatorname{End}_k(E)\bigr)
\]
is a closed immersion. Therefore \(G/N\) is affine.

\smallskip
\phantomsection\label{case:sga3-VIB-11.17-b}%
\noindent\textup{(b)}
Suppose now that \(G\) and \(N\) are reduced, without
assuming that \(G\) is connected. Put
\[
  N'=N\cap G^0.
\]
Then \(G^0/N'\) is affine by \textup{(a)}. Moreover,
\(NG^0\) is an invariant subgroup of \(G\), and
\(G/NG^0\), being a quotient of the finite constant group
\(G/G^0\)
\textup{(cf.\ \hyperref[prop:sga3-VIA-5.5.1]
{VI\(_A\), 5.5.1})}, is itself a finite constant group.
The scheme \(G/N\) is therefore the disjoint union of the
fibers of
\[
  G/N\longrightarrow G/NG^0.
\]
By \hyperref[par:sga3-VIA-5.3.3]{VI\(_A\), 5.3.3}, all
these fibers are isomorphic to \(NG^0/N\), hence to
\(G^0/N'\). Thus \(G/N\) is affine.

\smallskip
\phantomsection\label{case:sga3-VIB-11.17-c}%
\noindent\textup{(c)}
Suppose that \(G\) is reduced, with \(N\) arbitrary. The
morphism
\[
  G\times N\longrightarrow N,\qquad
  (g,h)\longmapsto ghg^{-1},
\]
induces a morphism
\[
  (G\times N)_{\mathrm{red}}\longrightarrow N_{\mathrm{red}}.
\]
Since \(G\) is reduced and \(k\) is algebraically closed,
\[
  (G\times N)_{\mathrm{red}}=G\times N_{\mathrm{red}},
\]
so \(N_{\mathrm{red}}\) is an invariant subgroup of \(G\).
\textup{(This fails when \(G\) is not reduced; see
\hyperref[sga3:VIA:0.2]{VI\(_A\), 0.2}.)}

Hence, by \textup{\SGARefLink{case:sga3-VIB-11.17-b}{(b)}},%
\begingroup
\renewcommand{\thefootnote}{\fnsymbol{footnote}}%
\footnote[1]{\label{note:sga3-VIB-11.17-case-c-source-reading}
Source-reading note: the French re-edition prints
``by \textup{(a)}'' here, and the Reinhold comparison
lineage repeats that reference. At this stage \(G\) need
not be connected; the applicable preceding case is
\textup{\SGARefLink{case:sga3-VIB-11.17-b}{(b)}}.}
\endgroup\ 
\[
  G'=G/N_{\mathrm{red}}
\]
is affine. By
\hyperref[prop:sga3-VIA-5.6.1]{VI\(_A\), 5.6.1},
\[
  N'=N/N_{\mathrm{red}}
\]
is a finite \(k\)-group. Hence, by
\hyperref[thm:sga3-V-4.1]{Theorem 4.1 of Expos\'e V},
\[
  G/N=G'/N'
\]
is affine.

\smallskip
\phantomsection\label{case:sga3-VIB-11.17-d}%
\noindent\textup{(d)}
For arbitrary \(G\) and \(N\), base change of the
equivalence relation
\[
  G\times N\rightrightarrows G
\]
along \(G_{\mathrm{red}}\to G\) gives
\phantomsection\label{form:sga3-VIB-11.17-reduced-relation}%
\[
  G_{\mathrm{red}}\times N'
    \rightrightarrows G_{\mathrm{red}},
  \qquad
  N'=N\cap G_{\mathrm{red}}.
\]
Since the underlying spaces are the same \textup{(and the
quotient is the quotient ringed space)}, the morphism
\[
  G_{\mathrm{red}}/N'\longrightarrow G/N
\]
is a homeomorphism. The quotient
\(G_{\mathrm{red}}/N'\) is reduced, because
\[
  p:G_{\mathrm{red}}\longrightarrow G_{\mathrm{red}}/N'
\]
is faithfully flat. It follows that
\[
  (G/N)_{\mathrm{red}}
    \simeq G_{\mathrm{red}}/N'.
\]
The scheme on the right is affine by \textup{\SGARefLink{case:sga3-VIB-11.17-c}{(c)}}. Since
\(G/N\) is of finite type over \(k\)
\textup{(cf.\
\hyperref[thm:sga3-VIA-3.3.2]{VI\(_A\), 3.3.2})},
EGA I, 5.1.10 implies that \(G/N\) is affine.

Finally, let \(k\) be arbitrary and let \(\overline{k}\)
be an algebraic closure. By
\hyperref[concl:sga3-VIB-9.2-v]{9.2\textup{(v)}}, one has
\[
  (G\otimes_k\overline{k})/
    (N\otimes_k\overline{k})
  \simeq
  (G/N)\otimes_k\overline{k}.
\]
The first scheme is affine, hence so is the second.
Therefore \(G/N\) is affine by \textup{(fpqc)} descent
\textup{(cf.\ EGA IV\(_2\), 2.7.1)}. This proves
\hyperref[thm:sga3-VIB-11.17]{Theorem 11.17} when \(G\) is of finite type. To pass to the
general case, we need the following lemma.%
\footnote{\label{note:sga3-VIB-133}Editor's note: we have
added this lemma, taken from \textup{[DG70]}, Section III.3,
7.1.}

\medskip
\phantomsection\label{lem:sga3-VIB-11.17.1}%
\noindent\textbf{Lemma 11.17.1.---}
\emph{Let \((C_i\to A_i)_{i\in I}\) be a filtered direct
system of ring homomorphisms, all faithfully flat. Then}
\[
  A=\varinjlim_i A_i
\]
\emph{is faithfully flat over}
\[
  C=\varinjlim_i C_i.
\]

\emph{Proof.}
By \textup{[BAC]}, \SGARefLink{sec:sga3-VIB-3}{Section I.3}, Proposition 9, a ring
homomorphism \(B\to B'\) is faithfully flat if and only if
it is injective and \(B'/B\) is a flat \(B\)-module. Since
each \(C_i\to A_i\) is faithfully flat, there are exact
sequences
\phantomsection\label{seq:sga3-VIB-11.17.1-stage}%
\[
  0\longrightarrow C_i\longrightarrow A_i
   \longrightarrow A_i/C_i\longrightarrow0,
\]
and \(A_i/C_i\) is a flat \(C_i\)-module. Thus
\[
  (A_i/C_i)\otimes_{C_i}C
\]
is a flat \(C\)-module. Since direct limits are exact and
commute with tensor products, we obtain an exact sequence
\phantomsection\label{seq:sga3-VIB-11.17.1-limit}%
\[
  0\longrightarrow C\longrightarrow A
   \longrightarrow A/C\longrightarrow0
\]
and an isomorphism
\phantomsection\label{form:sga3-VIB-11.17.1-limit-quotient}%
\[
  \varinjlim_i\bigl((A_i/C_i)\otimes_{C_i}C\bigr)
    =(A/C)\otimes_CC=A/C.
\]
It follows that \(A/C\) is a flat \(C\)-module. Hence
\(A\) is faithfully flat over \(C\).

We now return to the proof of \hyperref[thm:sga3-VIB-11.17]{Theorem 11.17} in the general
case, where \(G\) is not assumed to be of finite type. Put
\[
  \mathcal A(G)=A,
  \qquad
  \mathcal A(N)=B=A/J.
\]
By \hyperref[prop:sga3-VIB-11.13]{11.13}, \(A\) is the
direct limit of an increasing filtered family
\((A_i)_{i\in I}\) of finitely generated sub-Hopf
algebras. Thus
\[
  G=\varprojlim_iG_i,
  \qquad
  G_i=\operatorname{Spec}(A_i),
\]
where the \(G_i\) are affine algebraic \(k\)-groups. Let
\(\Delta\) and \(\tau\) denote, respectively, the
comultiplication and antipode of \(A\), and put
\[
  \Delta^2=(\Delta\otimes\mathrm{id}_A)\circ\Delta.
\]

For each \(i\),
\[
  B_i=A_i/(J\cap A_i)
\]
is a quotient Hopf algebra of \(A_i\), so
\[
  N_i=\operatorname{Spec}(B_i)
\]
is a closed subgroup of \(G_i\). Since \(N\) is invariant
in \(G\), the morphism
\[
  G\times N\longrightarrow G,\qquad
  (g,n)\longmapsto gng^{-1},
\]
factors through \(N\). Equivalently, the pair \((A,J)\)
satisfies
\phantomsection\label{form:sga3-VIB-11.17-conjugation-coideal}%
\[
  \bigl(m_{13}\circ(\Delta\otimes\tau)\circ\Delta\bigr)(J)
    \subset A\otimes_kJ,
\]
where
\[
  m_{13}(a_1\otimes a_2\otimes a_3)
    =a_1a_3\otimes a_2.
\]
It follows that \((A_i,A_i\cap J)\) satisfies the analogous
property, so \(N_i\) is invariant in \(G_i\). Moreover,
\[
  \varinjlim_iB_i=B,
  \qquad
  \varprojlim_iN_i=N.
\]

By the finite-type case already proved, each
\textup{(fpqc)} sheaf quotient \(G_i/N_i\) is represented
by an affine \(k\)-group
\[
  Q_i=\operatorname{Spec}(C_i).
\]
Put
\[
  C=\varinjlim_iC_i.
\]
The \(Q_i\) form a filtered projective system of affine
\(k\)-groups, and its projective limit
\[
  Q=\varprojlim_iQ_i
\]
is the \(k\)-group \(\operatorname{Spec}(C)\)
\textup{(cf.\ EGA IV\(_3\), 8.2.3)}. We have an exact
sequence of \(k\)-groups
\phantomsection\label{seq:sga3-VIB-11.17-Q}%
\[
  1\longrightarrow N\longrightarrow G\longrightarrow Q.
\]

We claim that \(Q\) represents the \textup{(fpqc)} sheaf
quotient of \(G\) by \(N\). It is enough to verify that
\[
  G\longrightarrow Q
\]
is a covering for the \textup{(fpqc)} topology
\textup{(cf.\ IV, \SGARefLink{sga3:iv:scholium:3:3:2:1}{3.3.2.1} and 5.1.7.1)}. Every morphism
\[
  G_i\longrightarrow Q_i
\]
is faithfully flat
\textup{(cf.\ \hyperref[concl:sga3-VIB-9.2-xi]
{9.2\textup{(xi)}})}; equivalently, \(A_i\) is faithfully
flat over \(C_i\). Since
\[
  A=\varinjlim_iA_i,
  \qquad
  C=\varinjlim_iC_i,
\]
\hyperref[lem:sga3-VIB-11.17.1]{Lemma 11.17.1} shows that \(A\) is faithfully flat over
\(C\). Thus \(G\to Q\) is faithfully flat. It is also
affine, hence quasi-compact, and is therefore an
\textup{(fpqc)} covering. This completes the proof of
\hyperref[thm:sga3-VIB-11.17]{Theorem 11.17}.

\medskip
\phantomsection\label{par:sga3-VIB-11.18}%
\noindent\textbf{11.18. Complements.---}%
\footnote{\label{note:sga3-VIB-134}Editor's note: we have
added this subsection.}
\hyperref[thm:sga3-VIB-11.17]{Theorem 11.17} and its proof also give the following
results, taken from \textup{[DG70]}, Section III.3.7. Let
\(k\) be a field. We begin with a lemma
\textup{(cf.\ \textup{[An73]}, 2.3.3.2)} that will be used
later \textup{(see \hyperref[prop:sga3-VIB-12.10]{12.10})}.

\medskip
\phantomsection\label{lem:sga3-VIB-11.18.1}%
\noindent\textbf{Lemma 11.18.1.---}
\emph{Let \(u:G\to G'\) be a morphism of \(k\)-groups and
put \(N=\operatorname{Ker}(u)\). Suppose that \(G'\) is
affine and that \(G\) is of finite type.}
\begin{enumerate}
\item[\textup{\SGAInlineRefTarget{lem:sga3-VIB-11.18.1:item:i}(i)}]\phantomsection
  \label{concl:sga3-VIB-11.18.1-i}%
  \emph{The morphism
  \(\tau:G/N\to G'\) is a closed immersion. In particular,
  \(G/N\) is affine.}

\item[\textup{(ii)}]\phantomsection
  \label{concl:sga3-VIB-11.18.1-ii}%
  \emph{If, moreover, the morphism
  \(u^\natural:\mathcal O(G')\to\mathcal O(G)\) is
  injective, then \(\tau\) is an isomorphism. Thus \(G'\)
  is of finite type and \(u\) is faithfully flat.}
\end{enumerate}

Indeed, by \hyperref[prop:sga3-VIB-11.13]{11.13},
\(G'\) is the projective limit of a filtered system of
affine algebraic \(k\)-groups \(G'_i\). Let
\[
  u_i:G\longrightarrow G'\longrightarrow G'_i
\]
be the composite, and let \(N_i\) be its kernel. The
\(N_i\) form a filtered decreasing system of closed
subgroups of \(G\), whose intersection is \(N\). Since
\(G\) is Noetherian, there is an index \(i\) such that
\[
  N=N_i.
\]
The groups \(G\) and \(G'_i\) are of finite type.
Therefore, by
\hyperref[thm:sga3-VIA-3.2]{VI\(_A\), 3.2} and
\hyperref[prop:sga3-VIA-5.4.1]{5.4.1}, the quotient
\(G/N\) is a \(k\)-group of finite type, and \(u_i\) is
the composite of the faithfully flat projection
\[
  p:G\longrightarrow G/N
\]
and a closed immersion
\[
  \tau_i:G/N\hookrightarrow G'_i.
\]

Consider the commutative diagram
\begin{center}
  \phantomsection\label{diag:sga3-VIB-11.18.1-closed-immersion}%
  \begin{tikzcd}[row sep=large,column sep=large]
    G \arrow[r,"u"] \arrow[d,"p"'] &
    G' \arrow[d,"q_i"] \\
    G/N \arrow[r,"\tau_i"'] \arrow[ur,"\tau"] &
    G'_i\mathpunct{.}
  \end{tikzcd}
\end{center}
Since
\[
  q_i\circ\tau=\tau_i
\]
is a closed immersion and \(q_i\) is separated
\textup{(\(G'\) is separated; cf.\
\hyperref[sga3:VIA:0.3]{VI\(_A\), 0.3})}, \(\tau\) is a
closed immersion
\textup{(cf.\ EGA I, 5.4.4)}. Hence \(G/N\) is affine and
\[
  \tau^\natural:\mathcal O(G')\longrightarrow
    \mathcal O(G/N)
\]
is surjective, proving \textup{\SGARefLink{lem:sga3-VIB-11.18.1:item:i}{(i)}}.

If \(u^\natural:\mathcal O(G')\to\mathcal O(G)\) is also
injective, then so is \(\tau^\natural\). It follows that
\(\tau^\natural\) is an isomorphism, and hence so is
\(\tau\), since \(G'\) and \(G/N\) are affine. This proves
\textup{\SGARefLink{concl:sga3-VIB-11.18.1-i}{(ii)}}.

\medskip
\phantomsection\label{thm:sga3-VIB-11.18.2}%
\noindent\textbf{Theorem 11.18.2.---}
\emph{Let \(G\) and \(G'\) be affine \(k\)-groups, with
algebras \(A\) and \(A'\), and let
\(u:G\to G'\) be a morphism of \(k\)-groups. Put
\(N=\operatorname{Ker}(u)\), and let
\(\varphi:A'\to A\) be the morphism induced by \(u\).}
\begin{enumerate}
\item[\textup{\SGAInlineRefTarget{thm:sga3-VIB-11.18.2:item:i}(i)}]\phantomsection
  \label{concl:sga3-VIB-11.18.2-i}%
  \emph{If \(\varphi\) is injective, then \(u\) is
  faithfully flat and identifies \(G'\) with \(G/N\).}

\item[\textup{\SGAInlineRefTarget{thm:sga3-VIB-11.18.2:item:ii}(ii)}]\phantomsection
  \label{concl:sga3-VIB-11.18.2-ii}%
  \emph{One has
  \(G/N=\operatorname{Spec}(B)\), where
  \(B=\varphi(A')\). Thus \(u\) is the composite of the
  faithfully flat morphism \(G\to G/N\), corresponding to
  the inclusion \(B\hookrightarrow A\), and the closed
  immersion \(G/N\hookrightarrow G'\), corresponding to
  the surjection \(A'\to B\). Moreover, \(N\) is defined
  in \(G\) by the ideal \(AB_+\), where \(B_+\) denotes
  the augmentation ideal of \(B\).}

\item[\textup{\SGAInlineRefTarget{thm:sga3-VIB-11.18.2:item:iii}(iii)}]\phantomsection
  \label{concl:sga3-VIB-11.18.2-iii}%
  \emph{In particular, if \(u\) is a monomorphism, it is
  a closed immersion.}
\end{enumerate}

\emph{Proof.}
\textup{(i)}
Suppose that \(\varphi\) is injective and identify \(A'\)
with a sub-Hopf algebra of \(A\). By
\hyperref[prop:sga3-VIB-11.13]{11.13}, \(A\) is the
filtered union of finitely generated sub-Hopf algebras
\[
  A_i=\mathcal O(G_i).
\]
Put
\[
  A'_i=A'\cap A_i,
  \qquad
  G'_i=\operatorname{Spec}(A'_i),
\]
and let \(N_i\) be the kernel of the morphism
\[
  G_i\longrightarrow G'_i
\]
induced by \(A'_i\hookrightarrow A_i\). By \hyperref[lem:sga3-VIB-11.18.1]{Lemma 11.18.1},
\[
  G_i/N_i\simeq G'_i.
\]
Thus, for each \(i\), there is an exact sequence
\phantomsection\label{seq:sga3-VIB-11.18.2-stage}%
\[
  1\longrightarrow N_i\longrightarrow G_i
   \xrightarrow{\ p_i\ }G'_i\longrightarrow1,
\]
where \(p_i\) is faithfully flat. Since
\[
  G=\varprojlim_iG_i,
  \qquad
  G'=\varprojlim_iG'_i,
\]
we obtain an exact sequence
\begingroup
\renewcommand{\thefootnote}{\fnsymbol{footnote}}%
\footnote[2]{\label{note:sga3-VIB-11.18.2-limit-target-source-reading}
Source-reading note: the French re-edition prints \(G/N\)
as the final term of this limiting exact sequence, and the
Reinhold comparison lineage repeats it. At this point the
argument is proving that \(G'\) represents \(G/N\); taking
the projective limits of the preceding sequences gives
\(G'\) as the final term.}
\endgroup
\phantomsection\label{seq:sga3-VIB-11.18.2-limit}%
\[
  1\longrightarrow N\longrightarrow G
   \xrightarrow{\ p\ }G',
\]
where \(N=\varprojlim_iN_i\). By \hyperref[lem:sga3-VIB-11.17.1]{Lemma 11.17.1}, \(p\) is
faithfully flat \textup{(and affine)}. Therefore \(G'\)
represents the \textup{(fpqc)} sheaf quotient of \(G\) by
\(N\), proving \textup{\SGARefLink{thm:sga3-VIB-11.18.2:item:i}{(i)}}.

In general, \(B=\varphi(A')\) is a sub-Hopf algebra of
\(A\). Let
\[
  H=\operatorname{Spec}(B),
\]
and let \(N'\) be the kernel of the morphism \(G\to H\)
induced by \(B\hookrightarrow A\). By \textup{\SGARefLink{thm:sga3-VIB-11.18.2:item:i}{(i)}}, \(H\)
is identified with \(G/N'\), and \(u\) is the composite
of the projection \(G\to G/N'\) and the closed immersion
\[
  G/N'\hookrightarrow G'
\]
induced by the surjection \(A'\to B\). Hence \(N'=N\).
Moreover, by
\hyperref[concl:sga3-VIB-9.2-ii]{9.2\textup{(ii)}}, the square
\begin{center}
  \phantomsection\label{diag:sga3-VIB-11.18.2-cartesian}%
  \begin{tikzcd}[row sep=large,column sep=large]
    N \arrow[r] \arrow[d] &
    G \arrow[d,"p"] \\
    \operatorname{Spec}(k) \arrow[r,"\varepsilon"'] &
    H
  \end{tikzcd}
\end{center}
is cartesian. Here \(\varepsilon\) is the unit section of
\(H\), corresponding to the augmentation \(B\to k\).\footnote{Source note: the re-edition prints \(G'\) in the lower
right corner and calls \(\varepsilon\) its unit. The augmentation
\(B\to k\) instead defines the unit of \(H=\operatorname{Spec}(B)=G/N\).}
It follows that \(N\) is defined in \(G\) by the ideal
\(AB_+\). This proves \textup{\SGARefLink{thm:sga3-VIB-11.18.2:item:ii}{(ii)}}, and \textup{\SGARefLink{thm:sga3-VIB-11.18.2:item:iii}{(iii)}}
follows.

\medskip
\phantomsection\label{rem:sga3-VIB-11.18.3}%
\noindent\textbf{Remark 11.18.3.---}
Let \(G\) be an affine \(k\)-group and \(N\) an invariant
\(k\)-subgroup. Since
\[
  p:G\longrightarrow G/N
\]
is faithfully flat and quasi-compact,
\(\mathcal O(G/N)\) is, by IV, \SGARefLink{sga3:iv:remark:3:3:3:2}{3.3.3.2}, the subalgebra of
\(\mathcal O(G)\) formed by the functions \(\phi\) that
are right \(N\)-invariant; that is,
\[
  \phi(gh)=\phi(g)
\]
for every \(k\)-scheme \(S\), every \(g\in G(S)\), and
every \(h\in N(S)\). Let \(J\) be the ideal of
\[
  A=\mathcal O(G)
\]
defining \(N\). The preceding condition is equivalent to
\phantomsection\label{form:sga3-VIB-11.18.3-invariants}%
\[
  \Delta(\phi)-\phi\otimes1
    \in\mathcal O(G)\otimes J,
\]
where \(\Delta\) is the comultiplication of \(A\).

The preceding theorem may therefore be reformulated in
terms of Hopf algebras as follows.

\medskip
\phantomsection\label{cor:sga3-VIB-11.18.4}%
\noindent\textbf{Corollary 11.18.4.---}
\emph{Let \(k\) be a field, let \(A\) be a commutative
\(k\)-Hopf algebra, and put \(G=\operatorname{Spec}(A)\).}
\begin{enumerate}
\item[\textup{\SGAInlineRefTarget{cor:sga3-VIB-11.18.4:item:i}(i)}]\phantomsection
  \label{concl:sga3-VIB-11.18.4-i}%
  \emph{If \(B\) is a sub-Hopf algebra of \(A\), then
  \(A\) is faithfully flat over \(B\).}

\item[\textup{(ii)}]\phantomsection
  \label{concl:sga3-VIB-11.18.4-ii}%
  \emph{The map}
  \[
    N\longmapsto\mathcal O(G/N)
  \]
  \emph{is a bijection from the set of invariant
  subgroups of \(G\) to the set of sub-Hopf algebras of
  \(A\). Its inverse is}
  \[
    B\longmapsto\operatorname{Spec}(A/AB_+).
  \]
  \emph{Moreover, if \(J\) is the ideal of \(A\) defining
  \(N\), then}
  \phantomsection\label{form:sga3-VIB-11.18.4-invariant-ring}%
  \[
    \mathcal O(G/N)
      =
    \{\,x\in A\mid
      \Delta(x)-x\otimes1\in A\otimes J\,\}.
  \]
\end{enumerate}
